\documentclass[a4paper,11pt]{article}

% ============================================================================
% TRABALHO FINAL - CIENCIA DE DADOS (TP2)
% UNESP - Faculdade de Ciencias - Bauru - Departamento de Computacao
%
% Tema 6 - Criminalidade urbana por regiao
% Recorte: Rio Grande do Sul, por MUNICIPIO (agrupado em mesorregioes do IBGE)
% ============================================================================

\usepackage[T1]{fontenc}
\usepackage[utf8]{inputenc}
\usepackage[brazil]{babel}
\usepackage{graphicx}
\usepackage{booktabs}
\usepackage{array}
\usepackage{enumitem}
\usepackage{xcolor}
\usepackage{amsmath}
\usepackage{hyperref}
\usepackage{indentfirst}
\usepackage[top=2.2cm,bottom=2.2cm,left=2.2cm,right=2.2cm]{geometry}

\hypersetup{
    hidelinks,
    pdftitle={Criminalidade urbana no Rio Grande do Sul -- Trabalho Final de Ciencia de Dados},
    pdfauthor={Mario Masao Mukai e Davi Ferreira de Souza}
}
\urlstyle{same}

\setlength{\parindent}{1.25cm}
\setlength{\parskip}{0pt}
\graphicspath{{../figuras/}{figuras/}}

% Gerado por src/tabelas_relatorio.py -- NAO EDITAR A MAO.
% Regenere com: python src/tabelas_relatorio.py

\newcommand{\linhasmacro}{%
\texttt{AMEACA\_VD} &  & 8 & 547.285 & 17,6 & Ameaca \\
\texttt{FRAUDE} &  & 27 & 536.481 & 17,2 & Estelionato \\
\texttt{OUTROS} &  & 156 & 405.237 & 13,0 & Outros crimes \\
\texttt{FURTO\_OUTROS} & $\bullet$ & 16 & 375.058 & 12,0 & Furto simples \\
\texttt{LESAO} & $\bullet$ & 8 & 278.911 & 8,9 & Lesao corporal \\
\texttt{TRANSITO} &  & 7 & 207.639 & 6,7 & Lesao corporal culposa direca. \\
\texttt{HONRA\_DISCRIM} &  & 12 & 139.989 & 4,5 & Injuria \\
\texttt{DROGAS} & $\bullet$ & 2 & 120.831 & 3,9 & Entorpecentes - trafico \\
\texttt{ORDEM\_PUBLICA} &  & 3 & 97.695 & 3,1 & Perturbacao do trabalho ou do. \\
\texttt{ROUBO\_RUA} & $\bullet$ & 6 & 81.441 & 2,6 & Roubo a pedestre \\
\texttt{DANO} &  & 6 & 79.659 & 2,6 & Dano \\
\texttt{FURTO\_VEICULO} & $\bullet$ & 3 & 68.682 & 2,2 & Furto em veiculo \\
\texttt{FURTO\_ARROMBAMENTO} & $\bullet$ & 6 & 54.627 & 1,8 & Furto/arrombamento de residen. \\
\texttt{ROUBO\_OUTROS} & $\bullet$ & 32 & 43.041 & 1,4 & Extorsao \\
\texttt{SEXUAL} &  & 5 & 35.037 & 1,1 & Estupro de vulneravel \\
\texttt{ARMAS} &  & 9 & 27.563 & 0,9 & Posse irregular de arma de fo. \\
\texttt{ROUBO\_VEICULO} & $\bullet$ & 4 & 12.468 & 0,4 & Roubo de veiculo \\
\texttt{CVLI} & $\bullet$ & 4 & 5.948 & 0,2 & Homicidio doloso \\
}

\newcommand{\linhasmeso}{%
Metropolitana de Porto Alegre & 98 & 4.831.544 & 7.995 & 5.760 & 0,89 & $<$0,001 \\
Centro Ocidental & 31 & 529.887 & 7.458 & 6.033 & 0,02 & 0,644 \\
Sudoeste & 19 & 698.692 & 6.709 & 6.278 & -0,32 & 0,147 \\
Nordeste & 54 & 1.134.876 & 6.513 & 4.959 & 0,97 & 0,005 \\
Noroeste & 216 & 1.981.313 & 6.499 & 5.001 & 0,68 & $<$0,001 \\
Sudeste & 25 & 891.593 & 6.243 & 5.465 & 0,27 & 0,946 \\
Centro Oriental & 54 & 815.060 & 5.934 & 4.368 & 0,48 & 0,028 \\
}

\title{\textbf{Onde e quando patrulhar?}\\[2mm]
\large Padrões espaciais e temporais da criminalidade no Rio Grande do Sul,\\
por município e mesorregião\\
\normalsize a partir dos microdados de ocorrências da SSP-RS}

\author{
Mário Masao Mukai -- RA: 241022321\\
Davi Ferreira de Souza -- RA: 241024676\\[2mm]
\small Trabalho Final -- Ciência de Dados\\
\small Universidade Estadual Paulista ``Júlio de Mesquita Filho'' (UNESP)\\
\small Faculdade de Ciências -- Bauru -- Departamento de Computação
}
\date{Bauru -- SP, \today}

\begin{document}
\maketitle

% ----------------------------------------------------------------------------
\begin{abstract}
\noindent
Uma secretaria estadual de segurança precisa distribuir um efetivo limitado
entre municípios, regiões e turnos, e costuma fazê-lo pelo volume bruto de
ocorrências, que apenas reflete o tamanho das cidades. Investigamos se há um
padrão espacial e temporal identificável na criminalidade do Rio Grande do Sul
que justifique uma alocação mais informada. Usamos os microdados de ocorrências
da SSP-RS (2022--2025, 3,1 milhões de registros) combinados com a população e as
mesorregiões do IBGE, com análise exploratória, teste qui-quadrado de
independência e teste de Kruskal-Wallis. Encontramos que o tipo de ocorrência
depende do turno de forma estatisticamente significativa, ainda que de magnitude
modesta (arrombamento de madrugada, furto de dia, roubo e lesão à noite), que o
turno de pico é a noite nas sete mesorregiões, seguida de perto pela tarde, e que
a taxa por habitante difere entre regiões. Recomendamos uma escala de policiamento uniforme no pico noturno e
diferenciada por tipo de crime e região, priorizando municípios pela taxa por
habitante e tratando o litoral como caso sazonal, com a ressalva de que o dado
mede ocorrência registrada e a população de referência é a residente.

\vspace{2mm}
\noindent\textbf{Palavras-chave:} ciência de dados; segurança pública;
análise exploratória; teste de hipóteses; taxa por habitante; mesorregião.
\end{abstract}

% ============================================================================
\section{Contexto}\label{sec:contexto}
% ============================================================================
A segurança pública no Brasil é atribuição dos estados, e cabe às secretarias
estaduais distribuir um contingente policial finito por um território extenso e
heterogêneo. No Rio Grande do Sul são 497 municípios, de capitais regionais a
cidades de poucos milhares de habitantes, com dinâmicas de criminalidade
distintas. Na ausência de evidência, essa distribuição tende a seguir o volume
absoluto de ocorrências, critério que favorece automaticamente as cidades
grandes, onde há mais pessoas e, portanto, mais registros, e que nada informa
sobre o \emph{risco} enfrentado pelo cidadão ou sobre \emph{quando} o crime
ocorre.

Padrões espaciais e temporais permitem outro critério. Se o crime se concentra
em certos turnos, a mesma equipe rende mais quando escalada nas horas certas. Se
o risco por habitante difere entre regiões, a prioridade deixa de ser apenas a
capital. Este trabalho investiga esses padrões a partir de dados públicos e
abertos, com o objetivo de subsidiar uma decisão de alocação.

% ============================================================================
\section{Pergunta central}\label{sec:pergunta}
% ============================================================================
O Tema~6 pergunta se os tipos de ocorrência têm padrão espacial ou temporal
identificável. Ajustamos o recorte para torná-lo acionável por um órgão
estadual: \emph{no Rio Grande do Sul, os tipos de ocorrência têm padrão espacial
(por município e mesorregião) e temporal (dia da semana, horário) identificável,
a ponto de justificar a priorização de recursos por região e turno?}

A escolha de analisar o estado inteiro por município, em vez de uma única
cidade por bairro, decorre de três razões: a fonte é estadual e cobre todos os
municípios; a agregação municipal permite normalizar por população e comparar
regiões de forma justa; e a decisão relevante para a SSP-RS é entre municípios e
regiões, não entre bairros de uma só cidade. A resposta não termina em
``existe padrão'', e sim em uma recomendação de alocação.

% ============================================================================
\section{Dados e metodologia}\label{sec:metodo}
% ============================================================================

\subsection{Fontes dos dados}
Foram usadas duas fontes públicas:
\begin{itemize}[nosep]
  \item \textbf{SSP-RS}, Secretaria da Segurança Pública do Rio Grande do Sul:
  microdados de ocorrências criminais (\url{https://ssp.rs.gov.br}), anos
  2022--2025, todo o estado, com granularidade de ocorrência individual. Colunas
  utilizadas: data, hora, grupo, tipo de enquadramento, município, local e
  quantidade de vítimas. Acesso em 16/09/2026.
  \item \textbf{IBGE}: população residente por município (Censo 2022, via API
  de Agregados) e a correspondência município\,$\rightarrow$\,mesorregião (API de
  Localidades), em \url{https://servicodados.ibge.gov.br}. Acesso em 16/09/2026.
\end{itemize}
Os dados são agregados e não contêm informação pessoal identificável.

\subsection{Limpeza e preparação}
O pipeline (\texttt{src/carga.py}, \texttt{src/ibge.py}, \texttt{src/limpeza.py})
executa: (1) remoção de duplicatas pela chave (Sequência, Enquadramento, Data);
(2) derivação de atributos temporais (ano, mês, dia da semana, indicador de fim
de semana e turno: madrugada, manhã, tarde, noite); (3) normalização do nome
do município e \emph{join} com a tabela do IBGE, com um pequeno dicionário de
correção para as poucas grafias divergentes (por exemplo, Sant'Ana do Livramento
e Doutor Maurício Cardoso), atingindo cobertura integral dos registros; e (4)
agrupamento dos tipos de enquadramento em macrocategorias (por exemplo,
\texttt{ROUBO\_RUA}, \texttt{FURTO\_OUTROS}, \texttt{LESAO}, \texttt{DROGAS}) e a
marcação de ``crimes de rua'': as categorias sensíveis a patrulhamento
ostensivo, por oposição a fraudes e crimes contra a honra, que são registrados na
delegacia mas não dependem de onde está a viatura. Os 314 tipos de enquadramento
foram reduzidos a 18 macrocategorias, das quais 9 compõem os crimes de rua,
33,4\% do total. Restaram 3.117.592 ocorrências ao fim da limpeza; a perda é
desprezível, porque os registros têm data e hora completas.

A Tabela~\ref{tab:macro} mostra o resultado dessa classificação. Duas
observações sobre ela. A primeira é que as três maiores categorias
(\texttt{AMEACA\_VD}, \texttt{FRAUDE} e \texttt{OUTROS}, juntas 47,8\% dos
registros) ficam \emph{fora} do recorte de crimes de rua, o que confirma a
necessidade do filtro: analisar o horário do estelionato ou da ameaça registrada
na delegacia diria pouco sobre onde posicionar a viatura. A segunda é que a
categoria \texttt{OUTROS} absorve 156 dos 314 tipos de enquadramento, mas apenas
13,0\% dos registros: metade do vocabulário da SSP-RS é cauda longa de baixa
frequência, e nenhum desses tipos isolados chega a 1\% do total.

\begin{table}[htbp]
\centering
\small
\caption{Macrocategorias, em ordem de frequência. O marcador $\bullet$ indica as
9 categorias classificadas como ``crime de rua'', sensíveis a patrulhamento
ostensivo, que somam 33,4\% dos registros.}
\label{tab:macro}
\begin{tabular}{l c r r r l}
\toprule
Macrocategoria & Rua & Tipos & Registros & \% & Enquadramento mais frequente \\
\midrule
\linhasmacro
\midrule
\textbf{Total} & & \textbf{314} & \textbf{3.117.592} & \textbf{100,0} & \\
\bottomrule
\end{tabular}
\end{table}

\subsection{Taxa por habitante}
A contagem bruta não é comparável entre municípios: o maior sempre lidera. Por
isso a comparação espacial usa a \textbf{taxa média anual por 100 mil
habitantes}: as ocorrências do período divididas pela população (Censo 2022) e
pelo número de anos do recorte, multiplicadas por 100\,000. A divisão pelo número
de anos mantém o numerador (somado sobre vários anos) na mesma escala do
denominador (população de um ano).

\subsection{Método de análise}
Além da análise exploratória, empregamos duas técnicas formais coerentes com a
natureza dos dados. O \textbf{qui-quadrado de independência} testa a associação
entre variáveis categóricas: aplicamo-lo a macrocategoria~$\times$~turno e a
macrocategoria~$\times$~mesorregião, com $H_0$ de independência, reportando
$\chi^2$, graus de liberdade, $p$-valor, o $V$ de Cramér (tamanho do efeito) e os
resíduos padronizados. O \textbf{teste de Kruskal-Wallis}, não-paramétrico,
verifica se a taxa por 100 mil habitantes dos municípios difere entre as
mesorregiões sem supor normalidade.

\subsection{Validação das premissas dos testes}\label{sec:premissas}
A escolha dos dois testes não é arbitrária, e as premissas foram verificadas
antes de aplicá-los.

Para o \textbf{qui-quadrado}, a condição de validade é que as frequências
esperadas sejam suficientemente grandes; a regra usual pede pelo menos 5 em cada
célula. Aqui a tabela macrocategoria~$\times$~turno tem 9 linhas e 4 colunas
sobre 1.041.007 ocorrências de crime de rua, e a menor frequência esperada fica
na ordem de centenas, muito acima do limite. O teste $G$ (razão de
verossimilhanças) seria uma alternativa equivalente e converge para o mesmo
resultado nesse regime de amostra; ficamos com o $\chi^2$ por ser o mais
difundido e por permitir a leitura direta dos resíduos padronizados, que é o que
de fato interpretamos.

Para a comparação da taxa entre mesorregiões, a escolha natural seria a ANOVA,
que porém exige normalidade dentro de cada grupo e variâncias homogêneas. As duas
premissas falham. O teste de Shapiro-Wilk rejeita a normalidade do conjunto das
497 taxas municipais ($W = 0{,}956$, $p \approx 4{,}9\times10^{-11}$), e a
distribuição é assimétrica à direita (assimetria $0{,}88$, curtose em excesso
$1{,}25$), como esperado de uma variável limitada em zero e com cauda longa de
municípios pequenos. Dentro dos grupos o quadro se repete: a normalidade é
rejeitada em 4 das 7 mesorregiões (Tabela~\ref{tab:meso}), justamente as de maior
número de municípios. O teste de Levene também rejeita a homogeneidade das
variâncias ($W = 3{,}28$, $p = 0{,}004$). Com as duas premissas violadas, o
Kruskal-Wallis, que compara distribuições por postos e não supõe nem normalidade
nem variâncias iguais, é a opção correta.

\begin{table}[htbp]
\centering
\small
\caption{Mesorregiões do Rio Grande do Sul. A taxa ponderada divide as
ocorrências da região pela sua população total; a mediana municipal trata cada
município como uma observação. As duas últimas colunas são a assimetria das taxas
municipais e o $p$-valor do teste de Shapiro-Wilk dentro da região.}
\label{tab:meso}
\begin{tabular}{l r r r r r r}
\toprule
 & & & \multicolumn{2}{c}{Taxa/100 mil hab.} & & \\
\cmidrule(lr){4-5}
Mesorregião & Mun. & População & Ponderada & Mediana & Assim. & $p$ (S-W) \\
\midrule
\linhasmeso
\bottomrule
\end{tabular}
\end{table}

\subsection{Alternativas consideradas}\label{sec:alternativas}
Três caminhos foram avaliados e descartados, por razões que vale registrar.

O primeiro foi a análise \textbf{por bairro}, que é o recorte sugerido no
enunciado do tema. Chegamos a processar o campo de bairro dos microdados e a
construir uma rotina de fusão de grafias por similaridade, cuja auditoria ficou
em \texttt{data/processed/auditoria\_fusao\_bairros.csv}: ela mostra casos como
\textsc{Jardom Botanico}~$\rightarrow$~\textsc{Jardim Botanico} e
\textsc{S Geraldo}~$\rightarrow$~\textsc{Sao Geraldo}. O problema é que o campo é
de preenchimento livre, sem padronização entre delegacias, e não existe
população por bairro no Censo que permita calcular taxa. Sem denominador, a
análise por bairro só produziria contagem bruta, que é exatamente o critério que
este trabalho questiona. Por isso o recorte subiu para o município.

O segundo foi uma \textbf{regressão de Poisson} da contagem de ocorrências por
município, com a população como \emph{offset}. Ela permitiria estimar o efeito de
covariáveis sobre a taxa, mas exigiria covariáveis municipais (renda, densidade,
estrutura etária) que estão fora do recorte de dados declarado, e responderia a
uma pergunta de \emph{explicação} que o Tema~6 não faz. A pergunta é se o padrão
existe e é acionável, não quais fatores o causam.

O terceiro foi o \textbf{agrupamento} (\emph{clustering}) dos municípios pelo
perfil horário, para descobrir tipologias de criminalidade. É atraente, mas
produziria grupos sem rótulo externo que os valide, e a decisão da SSP-RS já é
tomada sobre unidades administrativas existentes: as mesorregiões, que têm
comando regional. Um agrupamento estatístico que cortasse as mesorregiões ao meio
seria interessante e inútil para a alocação.

% ============================================================================
\section{Resultados}\label{sec:resultados}
% ============================================================================

\subsection{Quando o crime acontece}
A Figura~\ref{fig:heatmap} mostra a concentração de crimes de rua por dia e hora.
A atividade é baixa na madrugada (0h--6h), sobe ao longo do dia e atinge o
patamar mais alto entre a tarde e a noite. O fim de semana desloca parte das
ocorrências para a madrugada de sábado e domingo, quando as demais horas se
mantêm ativas.

\begin{figure}[htbp]
    \centering
    \includegraphics[width=0.92\textwidth]{f1_heatmap_dia_hora.pdf}
    \caption{Ocorrências de crime de rua por dia da semana e hora do dia
    (RS, 2022--2025).}
    \label{fig:heatmap}
\end{figure}

\subsection{Cada crime tem sua hora}
Agregar todos os crimes esconde que cada tipo tem seu horário. A
Figura~\ref{fig:perfil} normaliza a distribuição horária por macrocategoria: o
furto comum concentra-se de dia, ao passo que roubo e lesão sobem à noite. É essa
diferença de perfil que justifica uma escala de policiamento variável por turno,
e não apenas mais efetivo no horário de pico geral.

\begin{figure}[htbp]
    \centering
    \includegraphics[width=0.9\textwidth]{f2_perfil_horario_por_tipo.pdf}
    \caption{Distribuição horária por macrocategoria (crimes de rua): cada tipo
    tem um horário de pico distinto.}
    \label{fig:perfil}
\end{figure}

\subsection{Onde a taxa é maior}
Convertidas em taxa por habitante, as 15 maiores taxas municipais
(Figura~\ref{fig:municipios}) trazem oito municípios litorâneos, todos de pequeno
porte, e cidades de poucos milhares de habitantes como Esmeralda e Miraguaí.
Ambos os casos são artefatos que discutimos adiante. Ao restringir a municípios
com 50 mil habitantes ou mais (Figura~\ref{fig:grandes}), aparecem cidades do
interior como Vacaria e Cruz Alta ao lado de Porto Alegre: taxa elevada não é
exclusividade do litoral nem da capital.

\begin{figure}[htbp]
    \centering
    \includegraphics[width=0.82\textwidth]{f3_taxa_municipio.pdf}
    \caption{Taxa média anual por 100 mil habitantes, 15 maiores municípios
    (população: Censo 2022). O topo é dominado por municípios litorâneos.}
    \label{fig:municipios}
\end{figure}

\begin{figure}[htbp]
    \centering
    \includegraphics[width=0.82\textwidth]{f3b_taxa_municipio_grandes.pdf}
    \caption{Taxa média anual por 100 mil habitantes entre municípios de 50 mil
    habitantes ou mais, controlando o efeito das cidades pequenas.}
    \label{fig:grandes}
\end{figure}

\subsection{Intensidade por mesorregião}
Agregada por mesorregião e ponderada pela população
(Figura~\ref{fig:mesorregiao}), a taxa varia da região Metropolitana de Porto
Alegre (cerca de 8.000 ocorrências/ano por 100 mil hab.) ao Centro Oriental
(cerca de 5.900), uma diferença próxima de um terço.

Vale distinguir essa leitura da que o teste de Kruskal-Wallis faz, porque as duas
não medem a mesma coisa e a Tabela~\ref{tab:meso} mostra que elas discordam. A
taxa ponderada trata a região como uma unidade e é dominada pelos seus municípios
grandes: por ela, a Metropolitana lidera com folga, puxada por Porto Alegre e
pela conurbação. A mediana municipal trata cada município como uma observação, e
por ela a Metropolitana cai para terceira (5.760), atrás do Sudoeste (6.278) e do
Centro Ocidental (6.033). A explicação é que a Metropolitana reúne 98 municípios
muito desiguais, dos balneários do litoral norte a cidades-dormitório de taxa
baixa, e a mediana ignora o peso populacional que a ponderação captura.

As duas respostas são corretas para perguntas diferentes. Para dimensionar o
efetivo de uma região inteira, que é a decisão orçamentária, importa a taxa
ponderada, porque o efetivo atende pessoas. Para saber se \emph{o município
típico} de uma região enfrenta risco maior que o de outra, importa a mediana. O
Kruskal-Wallis responde à segunda pergunta, e é dela que trata a
Seção~\ref{sec:testes}.

\begin{figure}[htbp]
    \centering
    \includegraphics[width=0.82\textwidth]{f4_taxa_mesorregiao.pdf}
    \caption{Taxa média anual por 100 mil habitantes por mesorregião
    (ponderada pela população).}
    \label{fig:mesorregiao}
\end{figure}

O mapa (Figura~\ref{fig:mapa}) mostra a distribuição dessa taxa no território:
a região Metropolitana de Porto Alegre e o Centro Ocidental são as mais intensas,
e o Centro Oriental a menos. A malha geográfica foi obtida da API de malhas do
IBGE.

\begin{figure}[htbp]
    \centering
    \includegraphics[width=0.7\textwidth]{f8_mapa_mesorregioes.pdf}
    \caption{Taxa média anual por 100 mil habitantes por mesorregião do Rio
    Grande do Sul (mapa coroplético).}
    \label{fig:mapa}
\end{figure}

\subsection{O pico é a noite em todo o estado}
Cruzando espaço e tempo, a Figura~\ref{fig:turnoregiao} mostra a distribuição dos
crimes de rua por turno \emph{dentro} de cada mesorregião. A noite é o turno de
maior concentração nas sete regiões, o que simplifica a decisão, mas convém
registrar a margem: ela supera a tarde por 1 ponto percentual na Metropolitana
(31\% contra 30\%) e no Nordeste, e por 3 a 5 pontos nas demais. Noite e tarde
formam, na prática, um bloco contínuo de maior atividade. A diferença regional
clara está na madrugada, que pesa mais no interior (21\% das ocorrências no
Sudoeste, no Sudeste e no Centro Ocidental) do que na região Metropolitana
(15\%).

\begin{figure}[htbp]
    \centering
    \includegraphics[width=0.75\textwidth]{f7_turno_mesorregiao.pdf}
    \caption{Percentual de crimes de rua por turno dentro de cada mesorregião:
    o pico é a noite em todas.}
    \label{fig:turnoregiao}
\end{figure}

\subsection{Sazonalidade: o efeito do veraneio}
A Figura~\ref{fig:sazonalidade} contrasta a distribuição mensal das ocorrências no
litoral e no restante do estado: enquanto o interior é praticamente plano ao
longo do ano, o litoral salta para 12,5\% de suas ocorrências em janeiro (contra
8,6\% no resto), com novo pico em dezembro. O padrão coincide com a temporada de
veraneio, quando a população \emph{presente} é muito maior que a residente do
Censo, o que infla a taxa anual desses municípios.

\begin{figure}[htbp]
    \centering
    \includegraphics[width=0.9\textwidth]{f6_sazonalidade_litoral.pdf}
    \caption{Distribuição mensal das ocorrências: litoral \emph{versus} restante
    do estado. O litoral concentra ocorrências no verão.}
    \label{fig:sazonalidade}
\end{figure}

\subsection{Testes de hipótese}\label{sec:testes}
O qui-quadrado entre macrocategoria e turno rejeita a independência
($\chi^2 = 40{.}607$, $\mathrm{gl} = 24$, $p \approx 0$): o tipo de ocorrência
\emph{não} é independente do turno. Com 3,1 milhões de registros, porém, a
significância é praticamente garantida; o que importa é a magnitude, e o $V$ de
Cramér de $0{,}11$ indica uma associação real, porém modesta: o turno molda o
\emph{tipo} de crime, e responde por pouco da variação em \emph{se} há crime. Os
resíduos padronizados
(Figura~\ref{fig:residuos}) revelam o padrão: o furto por arrombamento concentra
na madrugada, o furto comum de manhã e à tarde, e roubo e lesão à noite; as
mesmas células mostram forte \emph{ausência} onde o crime não ocorre (o furto
comum praticamente some à noite).

O qui-quadrado entre macrocategoria e mesorregião também rejeita a independência
($\chi^2 = 50{.}368$, $\mathrm{gl} = 48$, $p \approx 0$; $V = 0{,}09$): o perfil
de crimes varia entre regiões, novamente com efeito modesto.

Por fim, o teste de Kruskal-Wallis compara a distribuição da taxa por 100 mil
habitantes dos 497 municípios, agrupados pela mesorregião a que pertencem, com
$H_0$ de que as sete distribuições são iguais. A hipótese é rejeitada
($H = 30{,}3$, $\mathrm{gl} = 6$, $p \approx 3{,}4\times10^{-5}$): o município
típico de uma região não enfrenta o mesmo risco que o de outra. O tamanho do
efeito, porém, acompanha o que os qui-quadrados já indicavam:
$\eta^2_H = 0{,}05$, ou seja, a mesorregião responde por cerca de 5\% da
variação nos postos das taxas municipais. Os outros 95\% são diferenças
\emph{dentro} de cada região, o que é coerente com a Tabela~\ref{tab:meso}: a
dispersão interna da Metropolitana, que vai de balneário a cidade-dormitório, é
maior que a distância entre as medianas regionais. A consequência prática é
direta, e sustenta a recomendação da Seção~\ref{sec:conclusao}: a região é um
critério de priorização fraco quando usada sozinha, e a taxa municipal é o
critério que de fato discrimina.

\begin{figure}[htbp]
    \centering
    \includegraphics[width=0.7\textwidth]{f5_residuos_chi2.pdf}
    \caption{Resíduos padronizados do $\chi^2$ (macrocategoria $\times$ turno):
    vermelho ocorre mais que o esperado sob independência; azul, menos.}
    \label{fig:residuos}
\end{figure}

% ============================================================================
\section{Limitações}\label{sec:limitacoes}
% ============================================================================
Quatro limitações qualificam as conclusões. Primeiro, o dado mede ocorrência
registrada, e não crime real: a subnotificação varia por tipo de crime e por
região, de modo que parte do padrão espacial pode refletir diferenças na
propensão a registrar. Segundo, a população de referência é a residente do Censo
2022, e não a população presente; como mostra a Figura~\ref{fig:sazonalidade}, os
municípios litorâneos aparecem com taxa inflada pela população flutuante do
verão, e sua taxa anual superestima o risco do morador fixo. Terceiro, os
municípios pequenos têm poucos casos e, portanto, taxas instáveis: Esmeralda, com
3.195 habitantes, ocupa a quinta posição da Figura~\ref{fig:municipios} com 1.459
ocorrências no período, e uma variação de poucas dezenas de registros a moveria
várias posições. Quarto, o recorte é o RS em 2022--2025, e as conclusões não se
estendem automaticamente a outros estados ou períodos.

% ============================================================================
\section{Conclusão e recomendação}\label{sec:conclusao}
% ============================================================================
A criminalidade no Rio Grande do Sul não se distribui ao acaso, nem no tempo nem
no espaço. No tempo, o tipo de ocorrência depende do turno de forma
significativa, ainda que modesta em magnitude, com um padrão nítido:
arrombamento de madrugada, furto de dia, roubo e lesão à noite. No espaço, a taxa
por habitante difere entre as mesorregiões. E o cruzamento das duas dimensões
simplifica a decisão: o pico é a noite em todas as regiões.

Portanto, a recomendação para a SSP-RS é uma escala de policiamento
\emph{uniforme quanto ao pico}, com reforço noturno em todo o estado, e
\emph{diferenciada quanto à composição}: patrulhamento ostensivo noturno contra
roubo e lesão nos centros urbanos, e vigilância de madrugada contra arrombamento,
com peso maior no interior, onde a madrugada é mais crítica. Na priorização entre
municípios, o critério deve ser a taxa por habitante (Figura~\ref{fig:grandes}), e
não o volume bruto. Essa priorização deve ser feita \emph{município a município},
e não por região: como a mesorregião responde por apenas 5\% da variação, tratar
uma região inteira como prioritária desperdiça efetivo nos seus municípios de taxa
baixa e ignora municípios críticos nas regiões de taxa média.

Com a ressalva de que o dado é de ocorrência registrada, sujeito a
subnotificação desigual; de que a taxa usa a população residente, o que faz o
litoral parecer de alto risco permanente quando seu problema é sazonal, a ser
tratado com reforço no verão em vez de prioridade o ano todo; e de que a
associação estatística encontrada, embora real, é de magnitude modesta, e deve
orientar a alocação sem substituir o conhecimento operacional de cada comando
regional.

% ============================================================================
% REFERENCIAS
% ============================================================================
\begin{thebibliography}{9}
\bibitem{sspRS}
Secretaria da Segurança Pública do Rio Grande do Sul.
\textit{Dados abertos --- Ocorrências criminais}.
Disponível em: \url{https://ssp.rs.gov.br}. Acesso em: 16 set. 2026.

\bibitem{ibgeCenso}
Instituto Brasileiro de Geografia e Estatística (IBGE).
\textit{Censo Demográfico 2022 --- População residente por município} (API de
Agregados) e \textit{API de Localidades} (mesorregiões).
Disponível em: \url{https://servicodados.ibge.gov.br}. Acesso em: 16 set. 2026.

\bibitem{scipy}
Virtanen, P. et al. \textit{SciPy 1.0: Fundamental Algorithms for Scientific
Computing in Python}. Nature Methods, v. 17, p. 261--272, 2020.
\end{thebibliography}

\end{document}
